\documentclass[11pt,aspectratio=169]{beamer}

\useinnertheme{metropolis}
\useoutertheme[progressbar=frametitle]{metropolis}
\usecolortheme{metropolis}
\usefonttheme{default}

\PassOptionsToPackage{round, authoryear}{natbib}
\usepackage{natbib}
\usepackage{appendixnumberbeamer}
\usepackage{amsmath, amssymb, amsthm}
\usepackage{bm}
\usepackage{booktabs}
\usepackage{xspace}
\usepackage{tikz}
\usepackage{graphicx}
\usetikzlibrary{arrows.meta, calc, positioning, shapes, backgrounds, fit}

\definecolor{semBlue}{RGB}{0,114,178}
\definecolor{semOrange}{RGB}{213,94,0}
\definecolor{semYellow}{RGB}{240,228,66}
\definecolor{semGreen}{RGB}{0,158,115}
\definecolor{semBg}{RGB}{255,253,236}
\definecolor{mDarkTeal}{HTML}{23373b}
\definecolor{blockTitleBg}{RGB}{225, 229, 233}
\definecolor{blockBodyBg}{RGB}{246, 247, 249}
\definecolor{citePurple}{RGB}{105, 55, 160}

\setbeamercolor{normal text}{fg=black,bg=semBg}
\setbeamercolor{frametitle}{fg=semBlue,bg=semBg}
\setbeamercolor{title}{fg=black,bg=semBg}
\setbeamercolor{progress bar}{fg=semOrange,bg=blockTitleBg}
\setbeamercolor{alerted text}{fg=semOrange}
\setbeamercolor{block title}{bg=blockTitleBg, fg=mDarkTeal}
\setbeamercolor{block body}{bg=blockBodyBg, fg=black}
\setbeamercolor{itemize item}{fg=semBlue}
\setbeamercolor{itemize subitem}{fg=semBlue}
\setbeamercolor{enumerate item}{fg=semBlue}
\setbeamertemplate{blocks}[rounded][shadow=false]
\setbeamertemplate{navigation symbols}{}
\setbeamerfont{author}{size=\large}
\setbeamertemplate{theorems}[numbered]

\hypersetup{
    colorlinks=true,
    citecolor=citePurple,
    linkcolor=citePurple,
    urlcolor=citePurple
}

\theoremstyle{plain}
\newtheorem{mytheorem}{Theorem}
\newtheorem{myproposition}{Proposition}
\newtheorem{mycorollary}{Corollary}

\theoremstyle{definition}
\newtheorem{mydefinition}{Definition}
\newtheorem{myexample}{Example}

\newenvironment{wideitemize}
  {\begin{itemize}\setlength{\itemsep}{0.72em}}
  {\end{itemize}}

\newenvironment{transitionframe}
  {\setbeamercolor{normal text}{fg=semBlue,bg=semBg}
   \setbeamercolor{frametitle}{fg=semBlue,bg=semBg}
   \begin{frame}[plain,c]}
  {\end{frame}
   \setbeamercolor{normal text}{fg=black,bg=semBg}
   \setbeamercolor{frametitle}{fg=semBlue,bg=semBg}}

\newcommand{\sectiontransition}[1]{
\begin{transitionframe}
\centering
{\Huge\bfseries #1}\\[0.8em]
{\color{semOrange}\rule{0.36\linewidth}{1.4pt}}
\end{transitionframe}
}

\newcommand{\backupbutton}[2]{\hyperlink{#1}{\beamergotobutton{#2}}}
\newcommand{\mainbutton}[2]{\hyperlink{#1}{\beamerreturnbutton{#2}}}

\newcommand{\R}{\mathbb{R}}
\newcommand{\I}{\mathbf{I}}
\newcommand{\G}{\mathbf{G}}
\newcommand{\q}{\mathbf{q}}
\newcommand{\x}{\mathbf{x}}
\newcommand{\y}{\mathbf{y}}
\newcommand{\0}{\mathbf{0}}
\newcommand{\SE}{\mathrm{SE}}
\newcommand{\PE}{\mathrm{PE}}
\newcommand{\E}{\mathrm{E}}

\title{\textbf{Strategic Elasticity}}
\subtitle{A relative measure of feedback in games}
\author{Zhigang Cao, \textbf{Sijie Wang}, Feng Zhu}
\institute{
Beijing Jiaotong University\\[0.35em]
\textit{Talk at HUST}\\[0.35em]
{\small 15-minute version}
}
\date{June 12, 2026}

\begin{document}

\maketitle

\begin{frame}{The puzzle}
\begin{wideitemize}
    \item Games are hard because players respond to one another.
    \item Weak feedback usually gives uniqueness, stability, and tractable comparative statics.
    \item The literature mostly measures feedback in \alert{absolute} terms:
    slopes, Jacobians, network effects, or spectral radii.
    \item But many economic restrictions are naturally \alert{relative}: bounded choices,
    positive baselines, and proportional changes.
\end{wideitemize}

\begin{block}{Question}
Can weak \alert{relative responsiveness} deliver good equilibrium properties
even when absolute responsiveness is large?
\end{block}
\end{frame}

\begin{frame}[label=main-absolute-relative]{Why absolute responsiveness can mislead}
\begin{columns}[T,onlytextwidth]
\begin{column}{0.48\textwidth}
\begin{block}{Standard view}
\begin{wideitemize}
    \item Strong network effects make uniqueness difficult.
    \item In linear games, conditions usually restrict $\G$.
    \item Individual characteristics $\q$ often disappear.
\end{wideitemize}
\end{block}
\end{column}
\begin{column}{0.48\textwidth}
\begin{block}{SEM view}
\begin{wideitemize}
    \item Bounds and stand-alone incentives reduce relative feedback.
    \item A large slope need not imply a large elasticity.
    \item $\G$, $\q$, and bounds jointly matter.
\end{wideitemize}
\end{block}
\end{column}
\end{columns}

\vspace{0.6em}
\centering
\backupbutton{backup-network-puzzle}{Network-game motivation}
\end{frame}

\sectiontransition{Idea}

\begin{frame}{Elasticity is log-log Lipschitz responsiveness}
For $\x,\y\in\R^M_{++}$, Thompson's metric is
\[
d_T(\x,\y)
=
\|\ln \x-\ln \y\|_\infty
=
\ln\max_i\left\{\frac{x_i}{y_i},\frac{y_i}{x_i}\right\}.
\]

For $f:K\subseteq\R^M_{++}\to\R_+$,
\[
\E_f^+(\x,\y)
=
\frac{d_T(f(\x),f(\y))}{d_T(\x,\y)}
=
\frac{|\ln f(\x)-\ln f(\y)|}
{\|\ln \x-\ln \y\|_\infty}.
\]

\begin{block}{Key observation}
Elasticity is a contraction modulus under log-log geometry.
\end{block}
\end{frame}

\begin{frame}{From functions to games}
\begin{wideitemize}
    \item A player's best response $br_i(\x_{-i})$ is the object that transmits strategic feedback.
    \item The strategic elasticity of player $i$ is the elasticity of $br_i$:
    \[
    \SE_i^+(\x_{-i})
    =
    \E_{br_i}^+(\x_{-i}).
    \]
    \item A game is \alert{inelastic} if every best response has point strategic elasticity below one.
    \item It is \alert{strongly inelastic} if the global strategic elasticity is below one.
\end{wideitemize}

\begin{block}{Interpretation}
A one-percent change in opponents' actions generates less than a one-percent
change in each player's best response.
\end{block}
\end{frame}

\sectiontransition{Results}

\begin{frame}{Fixed-point theorem: the engine}
\begin{mytheorem}[Fixed Points]
Let $K\subseteq\R^M_{++}$ be nonempty and log-convex, and let $F:K\to K$ be inelastic.
\begin{wideitemize}
    \item $F$ admits at most one fixed point.
    \item If a fixed point exists, iterating $F$ converges to it from every initial point.
    \item Existence follows from compact invariance, or from strong inelasticity plus closedness.
\end{wideitemize}
\end{mytheorem}

\begin{block}{Why this matters}
The method is a Banach-style contraction argument, but with Thompson's metric
instead of a normed metric.
\end{block}
\end{frame}

\begin{frame}{Equilibrium of inelastic games}
\begin{mytheorem}[Equilibrium]
For an inelastic game with log-convex strategy space $S\subseteq\R^N_{++}$:
\begin{wideitemize}
    \item there is at most one Nash equilibrium;
    \item if an equilibrium exists, best-response dynamics converge globally;
    \item existence is guaranteed under compactness, or under strong inelasticity and closedness.
\end{wideitemize}
\end{mytheorem}

\begin{block}{Take-away}
Inelasticity turns weak relative feedback into uniqueness, stability,
and dominance-solvability-type behavior.
\end{block}
\end{frame}

\begin{frame}{Comparative statics}
For a family of strongly inelastic games $\{\mathcal G(\bm{\theta})\}$,
let $\x^*(\bm{\theta})$ be the unique equilibrium.
\[
\E_{\x^*}^+(\bm{\theta})
\le
\max_{i\in N}
\frac{\PE_i^+(\bm{\theta})}
{1-\SE_i^+(\x^*(\bm{\theta}))}.
\]

\begin{columns}[T,onlytextwidth]
\begin{column}{0.47\textwidth}
\begin{block}{Direct effect}
\[
\PE_i^+(\bm{\theta})
\]
How parameters move player $i$'s best response directly.
\end{block}
\end{column}
\begin{column}{0.47\textwidth}
\begin{block}{Feedback multiplier}
\[
\frac{1}{1-\SE_i^+}
=1+\SE_i^++(\SE_i^+)^2+\cdots
\]
How strategic feedback amplifies the direct effect.
\end{block}
\end{column}
\end{columns}
\end{frame}

\sectiontransition{Application}

\begin{frame}{Linear best responses on bounded domains}
\begin{wideitemize}
    \item Consider bounded strategy spaces $[\mathbf a,\mathbf b]\subseteq\R^N_+$.
    \item Best response:
    \[
    br_i(\x_{-i})
    =
    \operatorname{med}
    \left\{
    a_i,\,
    q_i+\sum_{j\ne i}g_{ij}x_j,\,
    b_i
    \right\}.
    \]
    \item $\G=(g_{ij})$ measures interaction effects.
    \item $q_i$ is the autarkic optimum, or stand-alone incentive.
\end{wideitemize}

\begin{block}{What previous conditions miss}
Most uniqueness conditions focus on $\G$; SEM makes $\q$ and bounds visible.
\end{block}
\end{frame}

\begin{frame}{Strategic elasticity in the linear game}
\small
At interior points of player $i$'s best response,
\[
\SE_i^+(\x_{-i})
=
\frac{\sum_{j\ne i}|g_{ij}|x_j}
{q_i+\sum_{j\ne i}g_{ij}x_j}.
\]

\begin{block}{Reading the formula}
\begin{wideitemize}
    \item Numerator: total absolute influence from others.
    \item Denominator: the level of the best response.
    \item Bounds set elasticity to zero when the response is stuck at a boundary.
    \item Larger stand-alone incentives can dampen relative feedback.
\end{wideitemize}
\end{block}
\end{frame}

\begin{frame}{Main network-game implication}
\begin{myproposition}
For a bounded linear best-response game with $\q>\0$, the game is strongly
inelastic if, for every player $i$,
\[
q_i>
2\sum_{j:g_{ij}<0}|g_{ij}|b_j.
\]
\end{myproposition}

\begin{wideitemize}
    \item With pure strategic complements, $\G\ge0$, the condition reduces to $q_i>0$.
    \item With signed interactions, uniqueness depends on the balance between negative influence and stand-alone incentives.
    \item This answers a question about uniqueness conditions that depend jointly on $\G$ and $\q$.
\end{wideitemize}
\end{frame}

\begin{frame}{Relative pass-through}
\begin{myproposition}
If $\G\ge \0$ in the bounded linear network game, then
\[
\E_{\x^*}^+(\q)\le 1
\qquad
\text{for all }\q>\0.
\]
\end{myproposition}

\begin{block}{Interpretation}
A proportional increase in stand-alone incentives cannot generate a larger
proportional increase in equilibrium actions after network feedback.
\end{block}

\begin{wideitemize}
    \item Absolute pass-through can exceed one in network games.
    \item Relative pass-through is disciplined by strategic inelasticity.
\end{wideitemize}
\end{frame}

\sectiontransition{Message}

\begin{frame}{Contributions}
\begin{wideitemize}
    \item Introduce strategic elasticity as a relative measure of feedback in games.
    \item Show that inelastic mappings are contractions under Thompson's metric.
    \item Obtain uniqueness, existence, global stability, and quantitative comparative statics.
    \item Use SEM to explain why bounded network games can be well behaved even with large absolute network effects.
    \item Derive network-game conditions that incorporate both interaction structure $\G$ and individual characteristics $\q$.
\end{wideitemize}
\end{frame}

\begin{frame}{Take-away}
\begin{center}
{\Large
Strategic elasticity measures how much strategic feedback survives
after moving from absolute changes to relative changes.
}
\end{center}

\vspace{1em}

\begin{wideitemize}
    \item If feedback is relatively weak, equilibria are unique and stable.
    \item Thompson's metric is the geometry that makes this precise.
    \item In network games, strategic elasticity plays a role analogous to spectral radius,
    but it also incorporates scale, bounds, and stand-alone incentives.
\end{wideitemize}
\end{frame}

\appendix

\sectiontransition{Backup}

\begin{frame}[label=backup-network-puzzle]{Backup: network-game motivation}
\small
\begin{wideitemize}
    \item Strategic complementarity is a primary driver of equilibrium multiplicity
    \citep{cooper1988coordinating}.
    \item In network games, uniqueness usually requires weak network effects
    \citep{ballester2006s,bramoulle2014strategic}.
    \item \citet{belhaj2014network} show that bounded network games with strategic complements
    can have a unique equilibrium regardless of the strength of network effects.
    \item SEM explains the puzzle: bounded domains and positive baselines make the best-response
    map relatively inelastic.
\end{wideitemize}

\vspace{0.5em}
\mainbutton{main-absolute-relative}{Back}
\end{frame}

\begin{frame}{Backup: relation to existing fixed-point theorems}
\begin{wideitemize}
    \item Kennan's fixed point theorem is widely used in network economics.
    \item Du's fixed point theorem is useful for monotone dynamic economic models.
    \item The SEM fixed-point theorem subsumes both as special cases.
    \item The advantage is that SEM does not take monotonicity, concavity, or convexity
    as primitive assumptions.
\end{wideitemize}
\end{frame}

\begin{frame}{Backup: why $q_i$ matters in Thompson's metric}
In a normed metric, parallel translations do not change distance.

\[
\|(\x+\q)-(\y+\q)\|=\|\x-\y\|.
\]

In Thompson's metric, translations away from the origin can shrink relative distance:

\[
d_T(\x+\q,\y+\q)
=
\|\ln(\x+\q)-\ln(\y+\q)\|_\infty.
\]

\begin{block}{Economic reading}
Larger stand-alone incentives can reduce proportional sensitivity to peers.
This is why $\q$ appears in SEM but disappears from many norm-based conditions.
\end{block}
\end{frame}

\begin{frame}{Backup: references}
\footnotesize
\begin{wideitemize}
    \item Ballester, Calvo-Armengol, and Zenou (2006), ``Who's who in networks.''
    \item Belhaj, Bramoulle, and Deroian (2014), ``Network games under strategic complementarities.''
    \item Bulow, Geanakoplos, and Klemperer (1985), ``Multimarket oligopoly.''
    \item Cooper and John (1988), ``Coordinating coordination failures.''
    \item Du (1990), fixed point theorem for monotone mappings.
    \item Kennan (2001), uniqueness theorem for positive fixed points.
    \item Marshall (1890), \textit{Principles of Economics}.
    \item Thompson (1963), ``Certain contraction mappings in a partially ordered vector space.''
    \item Vives (1990, 2005), strategic complementarities and uniqueness.
\end{wideitemize}
\end{frame}

\end{document}
